% 3.2 倒格子和布里渊区
\section{倒格子和布里渊区}

若布拉维格子的基本平移矢量为 $\mathbf{t}_{1}$、$\mathbf{t}_{2}$、$\mathbf{t}_{3}$，则可以定义一组倒格子（reciprocal lattice）矢量 $\mathbf{g}_{1}$、$\mathbf{g}_{2}$、$\mathbf{g}_{3}$，满足
\begin{equation*}
\mathbf{g}_{i} \cdot \mathbf{t}_{j} = 2\pi\delta_{ij} \qquad (i, j = 1, 2, 3).
\tag{3.2.1}
\end{equation*}
这样定义的倒格子当然早已为 X 射线晶体学家所广泛使用。一些作者在式中定义 $\mathbf{g}_{i}$ 时不含因子 $2\pi$，因而在后面的理论中不得不使用“$2\pi$ 乘一个倒格子矢量”这种笨拙的说法。由式 (3.2.1)，$\mathbf{g}_{1}$、$\mathbf{g}_{2}$、$\mathbf{g}_{3}$ 可显式地写为
\begin{equation*}
\left\{
\begin{array}{l}
\mathbf{g}_{1} = \dfrac{2\pi(\mathbf{t}_{2} \times \mathbf{t}_{3})}{\mathbf{t}_{1} \cdot (\mathbf{t}_{2} \times \mathbf{t}_{3})}, \\[2ex]
\mathbf{g}_{2} = \dfrac{2\pi(\mathbf{t}_{3} \times \mathbf{t}_{1})}{\mathbf{t}_{2} \cdot (\mathbf{t}_{3} \times \mathbf{t}_{1})}, \\[2ex]
\mathbf{g}_{3} = \dfrac{2\pi(\mathbf{t}_{1} \times \mathbf{t}_{2})}{\mathbf{t}_{3} \cdot (\mathbf{t}_{1} \times \mathbf{t}_{2})}.
\end{array}
\right.
\tag{3.2.2}
\end{equation*}
若 $\mathbf{t}_{1}$、$\mathbf{t}_{2}$、$\mathbf{t}_{3}$ 是某一布拉维格子的基本矢量，则 $\mathbf{g}_{1}$、$\mathbf{g}_{2}$、$\mathbf{g}_{3}$ 通常是同一布拉维格子（在倒空间中）的基本矢量。但有四种情形，$\mathbf{g}_{1}$、$\mathbf{g}_{2}$、$\mathbf{g}_{3}$ 是同一晶系中另一种布拉维格子的基本矢量。例如，考虑面心立方格子，由表 3.1 其基本矢量为
\begin{equation*}
\left\{
\begin{array}{l}
\mathbf{t}_{1} = \frac{1}{2}a(0, 1, 1), \\
\mathbf{t}_{2} = \frac{1}{2}a(1, 0, 1), \\
\mathbf{t}_{3} = \frac{1}{2}a(1, 1, 0).
\end{array}
\right.
\tag{3.2.3}
\end{equation*}
把它们代入式 (3.2.2)，得到
\begin{equation*}
\left\{
\begin{array}{l}
\mathbf{g}_{1} = \frac{2\pi}{a}(-1, 1, 1), \\
\mathbf{g}_{2} = \frac{2\pi}{a}(1, -1, 1), \\
\mathbf{g}_{3} = \frac{2\pi}{a}(1, 1, -1),
\end{array}
\right.
\tag{3.2.4}
\end{equation*}
对照表 3.1 即可看出，这些矢量定义的是一个体心立方格子。因此，立方面心格子的倒格子是立方体心格子，反之亦然；正交面心格子的倒格子同样是正交体心格子，反之亦然。四方体心格子的倒格子是它自身，因为四方面心格子同时也是轴比不同的四方体心格子。与正空间中的布拉维格子一样，我们可以把倒格子任一点的位矢 $\mathbf{g}$ 用 $\mathbf{g}_{1}$、$\mathbf{g}_{2}$、$\mathbf{g}_{3}$ 表示；仿照式 (1.5.1)，
\begin{equation*}
\mathbf{g} = n_{1}\mathbf{g}_{1} + n_{2}\mathbf{g}_{2} + n_{3}\mathbf{g}_{3}
\tag{3.2.5}
\end{equation*}
其中 $n_{1}$、$n_{2}$、$n_{3}$ 为整数。

正空间坐标我们用 $x$、$y$、$z$，倒空间坐标用 $k_{x}$、$k_{y}$、$k_{z}$，因为倒空间中的一般矢量习惯上记作 $\mathbf{k}$。应当注意，如果把倒格子叠加在正格子上作图，那么 $\mathbf{g}_{1}$ 垂直于 $\mathbf{t}_{2}$ 和 $\mathbf{t}_{3}$，并不一定沿 $\mathbf{t}_{1}$ 的方向。在空间群的表示理论中究竟为什么要引入倒空间，到 3.4 节讨论平移群的表示理论时就会明白了。

表 3.3 列出 14 种布拉维格子的倒格子矢量。倒格子基矢 $\mathbf{g}_{1}$、$\mathbf{g}_{2}$、$\mathbf{g}_{3}$ 在所属晶系点群操作 $R$ 下的变换，既可以从倒空间中的图看出，也可以解析地计算：对式 (3.2.2) 两边施加 $R$ 即可，结果列于表 3.4。任一倒格子矢量 $\mathbf{g}$ 在 $R$ 下的变换，则可对式 (3.2.5) 两边施加 $R$ 并利用表 3.4 求出。

\begin{figure}[H]
\centering
\includegraphics[width=0.92\textwidth]{c3_t33.png}
\caption*{\textbf{表 3.3}\ 倒格子（原书扫描占位）。左起各列：布拉维格子、$\Gamma$ 记号、倒格子矢量 $\mathbf{g}_{1}, \mathbf{g}_{2}, \mathbf{g}_{3}$、$8\pi^{3}/V$。注：}
\end{figure}
\noindent（i）第 1 列列出正空间平移群 $\mathbf{T}$ 的布拉维格子。各布拉维格子的基本矢量见表 3.1。

\noindent（ii）第 2 列列出倒空间平移 $\mathbf{g}_{1}$、$\mathbf{g}_{2}$、$\mathbf{g}_{3}$ 相对于 $k_{x}$、$k_{y}$、$k_{z}$ 轴的坐标。例如 $\mathbf{g}_{1} = (p, q, r)$ 表示 $\mathbf{g}_{1} = p\mathbf{l} + q\mathbf{m} + r\mathbf{n}$，其中 $\mathbf{l}$、$\mathbf{m}$、$\mathbf{n}$ 是 $k_{x}$、$k_{y}$、$k_{z}$ 方向的单位矢量。

\noindent（iii）第 3 列给出 $(8\pi^{3}/V)$，即布里渊区的体积。

\begin{figure}[H]
\centering
\includegraphics[width=0.78\textwidth]{c3_t34_p1.png}
\caption*{\textbf{表 3.4}\ 算符作用在倒格子矢量上的结果（第一部分：单斜、正交；原书扫描占位）。}
\end{figure}
\begin{figure}[H]
\centering
\includegraphics[width=0.86\textwidth]{c3_t34_p2.png}
\caption*{\textbf{表 3.4}（第二部分：正交续、四方、三角与六角、立方；原书扫描占位）。}
\end{figure}
\begin{figure}[H]
\centering
\includegraphics[width=0.88\textwidth]{c3_t34_p3.png}
\caption*{\textbf{表 3.4}（第三部分：立方续；原书扫描占位）。注：}
\end{figure}
\noindent（i）为节省篇幅，表中只列出各晶系全对称点群（holosymmetric point group）诸元素的作用。由于 $I\mathbf{g} = -\mathbf{g}$ 对一切 $\mathbf{g}$ 成立，反演的作用也一并略去。

\noindent（ii）点群操作是主动的（active）。

\noindent（iii）倒格子的基本平移 $\mathbf{g}_{1}$、$\mathbf{g}_{2}$、$\mathbf{g}_{3}$ 与表 3.3 中的完全一致。

接下来我们讨论布里渊区（Brillouin 1930\textit{a}--\textit{c}, 1931, 1953）。\textit{第一布里渊区}（first Brillouin zone）本质上是晶体倒格子的一个晶胞；这是布里渊区要发挥任何有用作用的必要条件。惯常选取的特定晶胞是倒格子的 Wigner--Seitz 晶胞（见 1.5 节）（Jones 1960, Slater 1965\textit{b}）。这时布里渊区定义如下：

\noindent\textbf{定义 3.2.1}\quad 以倒格点之一为原点的全部 $\mathbf{k}$ 矢量的集合，若具有性质：从其中任一矢量出发，加上倒格子的平移矢量，不能到达更短长度的矢量，则这个集合构成所谓\textit{第一布里渊区}。

实际上，对三斜和单斜空间群我们将不使用这个定义（见下面的定义 3.2.2）。

按这个定义构造第一布里渊区时，我们选定某个倒格点为原点，作出由它连向其余倒格点的矢量 $\mathbf{g}$。若画出这些倒格子矢量 $\mathbf{g}$ 的垂直平分面，那么第一布里渊区就是围绕原点、由这些平面围成的最小体积。$\mathbf{g}$ 的垂直平分面的方程为
\begin{equation*}
\mathbf{k} \cdot \mathbf{g} = \tfrac{1}{2}|\mathbf{g}|^{2}.
\tag{3.2.6}
\end{equation*}
于是第一布里渊区的每一个面都由一个倒格子矢量 $\mathbf{g}$ 表征，并且当然就是 $\mathbf{g}$ 的垂直平分面。之所以在倒空间选取 Wigner--Seitz 晶胞而不是别的晶胞，是因为这个晶胞展示出倒格子的点群对称性（例如见 Delaunay (1932)）。然而，对属于低对称晶系的晶体，用式 (3.2.6) 作图构造布里渊区极其繁琐。

另一种值得在此一并考虑的晶胞是初基晶胞（见 1.5 节），即以 $\mathbf{k} = \mathbf{0}$ 为中心、棱与倒格子基本矢量 $\mathbf{g}_{1}$、$\mathbf{g}_{2}$、$\mathbf{g}_{3}$ 平行且等长的平行六面体。第一布里渊区的体积便由
\begin{equation*}
\mathbf{g}_{1} \cdot (\mathbf{g}_{2} \times \mathbf{g}_{3}) = \frac{8\pi^{3}}{V}
\tag{3.2.7}
\end{equation*}
给出，其中 $V$ 是相应正格子初基晶胞的体积；每种倒格子的 $\mathbf{g}_{1} \cdot (\mathbf{g}_{2} \times \mathbf{g}_{3})$ 值见表 3.3。对某些倒格子（正交 $P$、四方 $P$ 和立方 $P$ 空间群的倒格子），用倒空间初基晶胞定义的布里渊区与用 Wigner--Seitz 晶胞定义的惯用布里渊区完全相同。对另一些倒格子（正交 $I$、四方 $I$、六角 $P$、立方 $F$ 和立方 $I$ 空间群的倒格子），初基晶胞不能在与正空间格子相同的取向上展示布拉维格子的点群对称性，此时初基晶胞便不能用来定义布里渊区，见表 3.5 第 2 列。这八种布拉维格子几乎涵盖了实际研究过布里渊区的全部晶体，因此人们通常把选取倒空间的 Wigner--Seitz 晶胞当作第一布里渊区定义的一个必要部分，这并不奇怪。然而，对其余六种布拉维格子（三斜 $P$、单斜 $P$、单斜 $B$、正交 $C$、正交 $F$ 和三方 $R$ 空间群的），可以直接看出倒格子的初基晶胞也以正确的取向具有晶体布拉维格子的点群对称性。不过，只要可能，最好让布里渊区的界面平行于晶体中的对称面。这意味着对正交 $C$ 与 $F$ 空间群，在定义布里渊区时宜保留 Wigner--Seitz 胞，见图 3.6 与图 3.8。至于三方 $R$ 格子，似乎也没有什么特别有力的理由要用倒空间的初基晶胞去取代惯用的布里渊区定义。但对剩下的三种格子（三斜 $P$、单斜 $P$ 和单斜 $B$ 空间群的），如前所述，无论要画出布里渊区的形状，还是画出之后要在头脑中想象它，式 (3.2.6) 在实际应用中都十分困难，其形状细节在很大程度上依赖于所论晶体晶胞参数的具体取值。对单斜晶体，一些作者（Koster 1957, Sushkevich 1966）对每种布拉维格子各画出一个特定例子，Luehrmann (1968) 对单斜 $B$ 格子给出了两个例子，而 Kovalev (1965) 的详表则完全没有布里渊区图。因此，对单斜或三斜晶体采用初基晶胞定义布里渊区有一个优点：其总体形状（即棱数和面数）不依赖于晶体晶胞的轴比和轴间角。

\begin{table}[H]
\centering
\caption*{\textbf{表 3.5}\ 第 2 列标出：由倒空间初基晶胞定义的布里渊区是否以正确取向具有晶体的点群对称性；第 3 列标出：满足前一条件、且倒空间的 Wigner--Seitz 胞与初基胞给出相同布里渊区形状的那些布拉维格子。}
\begin{tabular}{@{}lcccclccc@{}}
\toprule
布拉维格子 & 2 & 3 & \phantom{xxxx} & 布拉维格子 & 2 & 3 \\
\midrule
三斜 $P$   & $\checkmark$ & $\times$ & & 四方 $P$ & $\checkmark$ & $\checkmark$ \\
单斜 $P$   & $\checkmark$ & $\times$ & & 四方 $I$ & $\times$ & --- \\
单斜 $B$   & $\checkmark$ & $\times$ & & 三方 $R$ & $\checkmark$ & $\times$ \\
正交 $P$   & $\checkmark$ & $\checkmark$ & & 六角 $P$ & $\times$ & --- \\
正交 $C$   & $\checkmark$ & $\times$ & & 立方 $P$ & $\checkmark$ & $\checkmark$ \\
正交 $I$   & $\times$ & --- & & 立方 $F$ & $\times$ & --- \\
正交 $F$   & $\checkmark$ & $\times$ & & 立方 $I$ & $\times$ & --- \\
\bottomrule
\end{tabular}
\end{table}

因此我们建议：大多数晶体的布里渊区仍用倒空间的 Wigner--Seitz 胞来定义，但对单斜和三斜空间群，应当改用倒空间的初基晶胞（Bradley and Cracknell 1970）。在构造图 3.2--3.4 时我们遵循了这一做法。

\noindent\textbf{定义 3.2.2}\quad 对单斜或三斜晶体，第一布里渊区取为倒格子的初基晶胞，即以 $\mathbf{k} = \mathbf{0}$ 为中心、棱与倒格子基本矢量 $\mathbf{g}_{1}$、$\mathbf{g}_{2}$、$\mathbf{g}_{3}$ 平行且等长的平行六面体。

最后还应指出一点。有人认为式 (3.2.6) 与电子、声子或磁振子的能量 $E(\mathbf{k})$ 随 $\mathbf{k}$ 变化曲线中出现不连续有特殊关联。其实并非如此，因为无论在扩展区图式还是重复区图式中，$E(\mathbf{k})$ 在倒空间处处都是 $\mathbf{k}$ 的准连续函数（Koopmans 定理）。划定区界时出现的不连续，不过是不同的 $E(\mathbf{k})$ 曲线在此分开。当然，有时可以细心地选取区界，使得 $E(\mathbf{k})$ 在区界处呈现特殊的简并；然而这一性质来源于“布里渊区应当展示布拉维格子的正确点群对称性”这一要求——我们并没有放弃这个要求——而不是使用 Wigner--Seitz 胞和式 (3.2.6) 的结果。

用式 (3.2.6) 构造第一布里渊区时，倒格子矢量 $\mathbf{g}$ 是由原点连向它的第一近邻、可能还有第二或更远近邻的矢量。还可以利用式 (3.2.6) 的更多平面，在第一布里渊区之外围出一个满足下列三个条件的体积，从而定义第二布里渊区：(i) 第二布里渊区的体积等于第一布里渊区的体积；(ii) 对第一布里渊区中的每个 $\mathbf{k}$ 矢量 $\mathbf{k}_{1}$，在第二布里渊区中恰有一个 $\mathbf{k}$ 矢量 $\mathbf{k}_{2}$，使得
\begin{equation*}
\mathbf{k}_{2} - \mathbf{k}_{1} = \mathbf{g},
\tag{3.2.8}
\end{equation*}
其中 $\mathbf{g}$ 是一个倒格子矢量，但对所有 $\mathbf{k}_{1}$ 不必相同。于是，第二布里渊区中的每一点都在某种意义上（3.3 节将更严格地定义）等价于第一布里渊区中的某一点；(iii) 第二布里渊区中的 $\mathbf{k}$ 矢量必须取满足条件 (i) 与 (ii) 的最短者。类似地可以定义第三、第四以及更高阶的布里渊区。

在给任一给定晶体定义布里渊区时，务必小心使用描述晶体\textit{全部}平移对称性的布拉维格子，而不要用晶胞更大的其他格子（Barron and Fischer 1959）。一旦完整认定了布拉维格子，倒格子与布里渊区的结构便唯一确定；它们与晶胞内部原子的具体排布完全无关。如上定义的布里渊区具有如下性质：

\noindent(i) 晶体的布拉维格子每个初基晶胞在区中对应一个态（对电子，泡利不相容原理允许自旋为 $+\frac{1}{2}$ 和 $-\frac{1}{2}$ 的两个电子占据每个这样的态；对声子或磁振子，当然不存在这样的不相容原理）（见式 (3.4.5)）。

\noindent(ii) 晶体中粒子或准粒子的能量在每个区内部处处连续。

\noindent(iii) 更高阶区的各分区可以通过适当的倒格子矢量平移而搬入第一区，这样每个区的态便在第一区上构成一条连续能带；这称为\textit{约化区图式}（reduced zone scheme）。

\noindent(iv) 具有第一布里渊区形状的胞可以相互拼合，完全填满倒空间，即相邻胞之间不留空隙。这称为\textit{重复区图式}（repeated zone scheme），常常有用。

在倒空间中定义区的另一种方法（Jones 1934\textit{a}）主张：只有那些对应的 Bragg 反射具有非零结构因子（structure factor）的平面才是真正的区界。这些 \textit{Jones 区}（Jones zones）在许多标准教科书中已有讨论（例如 Jones (1960), Kittel (1956), Mott and Jones (1936), Wilson (1953)），但如今似乎已很少使用。

\begin{figure}[H]
\centering
\includegraphics[width=0.92\textwidth]{fig3_1.pdf}
\caption{二维正方布拉维格子 $p$ 的前三个布里渊区（依原书重绘）。}
\end{figure}

我们用一个二维的例子来说明第一、第二及更高阶布里渊区的概念。二维正方布拉维格子 $p$ 的基本矢量可取为
\begin{equation*}
\left\{
\begin{array}{l}
\mathbf{t}_{1} = a\mathbf{i} \\
\mathbf{t}_{2} = a\mathbf{j}
\end{array}
\right.
\tag{3.2.9}
\end{equation*}
其中 $\mathbf{i}$、$\mathbf{j}$ 是 $x$、$y$ 方向的单位矢量。满足式 (3.2.1) 的倒格子矢量 $\mathbf{g}_{1}$、$\mathbf{g}_{2}$ 于是为
\begin{equation*}
\left\{
\begin{array}{l}
\mathbf{g}_{1} = \frac{2\pi}{a}\mathbf{j} \\
\mathbf{g}_{2} = \frac{2\pi}{a}\mathbf{i}.
\end{array}
\right.
\tag{3.2.10}
\end{equation*}
第一布里渊区是 $k_{x} = \pm\frac{1}{2}|\mathbf{g}_{2}| = \pm\pi/a$ 与 $k_{y} = \pm\frac{1}{2}|\mathbf{g}_{1}| = \pm\pi/a$ 之间的区域，也就是说，代入式 (3.2.6) 的矢量 $\mathbf{g}$ 是 $+\mathbf{g}_{1}$、$-\mathbf{g}_{1}$、$+\mathbf{g}_{2}$ 和 $-\mathbf{g}_{2}$。于是第一布里渊区就是图 3.1 中画有斜线阴影的正方形。对第二布里渊区，代入式 (3.2.6) 的矢量是 $\mathbf{g}_{1}+\mathbf{g}_{2}$、$-\mathbf{g}_{1}+\mathbf{g}_{2}$、$\mathbf{g}_{1}-\mathbf{g}_{2}$ 和 $-\mathbf{g}_{1}-\mathbf{g}_{2}$，第二布里渊区因而是图 3.1 中标有点点的区域。事实上我们只使用第一布里渊区，所以从现在起省略“第一”这个定语。

按照 Bouckaert, Smoluchowski, and Wigner (1936) 的说法，这些区的存在几乎同时被多位作者注意到（Bloch 1928, Morse 1930, Peierls 1930, Strutt 1928\textit{a},\textit{b}），而 Brillouin（1930\textit{a}--\textit{c}, 1931）指出了它们与 X 射线反射的联系。关于布里渊区已有一本教科书（Jones 1960）。
